\chapter{Algorithm 2: Deflected Subgradient Method}
\label{chap:dsm}

In this chapter, we apply the deflected subgradient method with Polyak step and with target level (SGPTL) to the ELM LASSO problem.

% ------------------------------------------------------------------
\section{Motivation: why the plain subgradient method is inadequate for ELM LASSO}
\label{sec:dsm_motivation}

The ELM LASSO problem
\begin{equation}\label{eq:dsm_objective}
  f(\mathbf{w}) = \frac{1}{2}\norm{\mathbf{X}\mathbf{w} - \mathbf{y}}_2^2 + \lambda_{\text{LASSO}}\norm{\mathbf{w}}_1,
\end{equation}
is nonsmooth due to the $\ell_1$ penalty. While the plain subgradient method guarantees convergence for convex nonsmooth problems, two structural limitations make it impractical for this problem:

\begin{enumerate}
\item \textbf{Zig-zagging and slow practical convergence:}
\\The subgradient update $\mathbf{w}_{i+1} = \mathbf{w}_i - \alpha_i \mathbf{g}_i$ uses only the current subgradient with no memory of previous iterations \cite{fundamental-relationships-sg}, which on poorly conditioned $\mathbf{X}^T\mathbf{X}$ produces severe oscillations. The complexity lower bound for first-order methods on nonsmooth convex problems is $\Omega(L^2/\varepsilon^2)$ \cite{nonsmooth-lower-bounds-frangioni,bubeck2015convex}, with $L$ a uniform bound on the subgradients. Plain subgradient already attains this optimal rate, but on ELM LASSO the constant $L$ depends on $\|\mathbf{X}\|_2$ and on the radius of the sublevel set containing the iterates, both of which can be large for poorly scaled or ill-conditioned data, so practical convergence is slow.

\item \textbf{Stringent stepsize requirements:}
\\Convergence requires the Decreasing Stepsize Scheduling condition \cite{fundamental-relationships-sg}, $\sum_i \alpha_i = \infty$ and $\sum_i \alpha_i^2 < \infty$, which forces the steps to shrink quickly (e.g.\ $\alpha_i = c/i$) and stalls practical progress on the accuracy targets we care about.
\end{enumerate}

Deflection addresses both issues by injecting memory into the search direction, and is the variant we adopt for ELM LASSO.

% ------------------------------------------------------------------
\subsection{The deflection mechanism}
\label{sec:dsm_deflection}

Deflected subgradient methods \cite{deflected-subgradient-method} replace the raw subgradient with a direction that blends the current subgradient with the previous search direction:
\begin{equation}\label{eq:deflected_dir}
  \mathbf{d}_i = \gamma_i\, \mathbf{g}_i + (1 - \gamma_i)\, \mathbf{d}_{i-1} \quad (i \ge 1), \qquad \mathbf{d}_0 = \mathbf{g}_0,
\end{equation}
where $\gamma_i \in (0,1]$ is the deflection parameter and the first iteration is handled separately (no previous direction is available). The update becomes $\mathbf{w}_{i+1} = \mathbf{w}_i - \alpha_i\, \mathbf{d}_i$. Setting $\gamma_i = 1$ at every step recovers the plain subgradient method, whereas $\gamma_i < 1$ gives weight to the memory term $\mathbf{d}_{i-1}$, smoothing out the trajectory. This damps the immediate direction reversals that make the plain subgradient method zig-zag on badly conditioned ELM LASSO instances. We pick $\gamma_i$ greedily to minimize $\norm{\mathbf{d}_i}$, which, coupled with the Polyak-style stepsize introduced below, maximizes the reduction achievable at iteration $i$ \cite{polyak-stepsize,deflected-subgradient-method}.

% ------------------------------------------------------------------
\section{Convergence framework: balancing deflection and stepsize}
\label{sec:dsm_convergence_conditions}

For deflected methods on ELM LASSO, the deflection $\gamma_i$ and stepsize $\alpha_i$ must be carefully coordinated. The theoretical framework is in \cite{deflected-subgradient-method}. We use the stepsize-restricted approach \cite{deflected-subgradient-method}:

\begin{equation}\label{eq:sr_step}
\alpha_i = \frac{\beta_i\,(f(\mathbf{w}_i) - f^*)}{\norm{\mathbf{d}_i}_2^2}, \quad \beta_i \leq \gamma_i.
\end{equation}

The condition $\beta_i \leq \gamma_i$ ensures stronger deflection is balanced by smaller steps, preventing large moves in memory-weighted non-descent directions. Since $f^*$ is unknown for ELM LASSO, we replace it with target level $f_{\mathrm{ref}} - \delta$ (Section~\ref{sec:dsm_target_level}), providing an adaptive estimate throughout the run.

\paragraph{Practical handling of the multiplier $\beta_i$.}\mbox{} \\
We take $\beta_i = \min(\beta,\gamma_i)$ with $\beta = 1$, the stepsize-restricted rule of \cite[Algorithm~SGPTL]{target-level-mechanism}; this enforces $\beta_i \le \gamma_i$ by construction, the hypothesis under which the deflected-Polyak rate (used in Theorem~\ref{thm:dsm_convergence}) is established in \cite{deflected-subgradient-method}, with original reference \cite{dantonio2009deflected}. The clip is operational on ELM~LASSO only because we pair it with $\gamma_i \ge \gamma_{\min} > 0$ on the deflection parameter, motivated below.

\subparagraph{Why $\gamma_{\min} > 0$ is required.}
Without a floor the closed-form greedy minimiser of
eq.~\eqref{eq:gamma_opt} can be projected onto $\gamma = 0$. This is
not a pathological corner case but the typical regime any descent run
enters within a few tens of iterations: once $\mathbf{w}_i$ approaches
a near-stationary region of $f$, consecutive subgradients
$\mathbf{g}_i, \mathbf{g}_{i-1}$ become well aligned, the parabola
coefficients of eq.~\eqref{eq:gamma_parabola} satisfy
$\norm{\mathbf{d}_{i-1}}_2^2 - \inner{\mathbf{g}_i}{\mathbf{d}_{i-1}}
\to 0^+$ from below the denominator $\norm{\mathbf{g}_i-\mathbf{d}_{i-1}}_2^2$, and the
unconstrained minimiser falls below $0$ --- it gets clipped to the
left boundary. The deflected direction collapses to $\mathbf{d}_i = \mathbf{d}_{i-1}$
(pure memory), $\gamma_i = 0$, and through the clip $\beta_i = 0$ and
the Polyak step $\alpha_i = 0$. The iterate freezes.

We therefore project the greedy minimiser of
eq.~\eqref{eq:gamma_opt} onto $[\gamma_{\min}, 1]$ instead of
$[0, 1]$, with $\gamma_{\min} = 0.05$ as the default. The choice of
$\gamma_{\min}$ trades two competing concerns. On one hand,
$\gamma_{\min}$ should be small enough that the floor rarely binds in
the pre-stationary phase, where the unconstrained
minimiser already lies inside $[0, 1]$: this preserves the
$\norm{\mathbf{d}_i}$-minimisation property that justifies the greedy
choice. On the other hand, $\gamma_{\min}$ should be large enough that
$\beta_i = \min(\beta, \gamma_i) \ge \gamma_{\min}$ keeps the Polyak
step $\alpha_i$ bounded away from zero and the iterate moving. We
sweep $\gamma_{\min} \in \{0.05, 0.1, 0.2, 0.5\}$ on the convergence
instance of Section~\ref{sec:convergence} and find final record gaps
in the $\sim 10^{-7}$--$10^{-5}$ range across the entire grid, with
$\gamma_{\min} = 0.05$ marginally best ($\bar f^{i} - f^{*} \approx 5\cdot 10^{-8}$
at $\rho = 0.7$); the algorithm is robust to the floor inside this
range and we adopt $\gamma_{\min} = 0.05$ in the rest of the report.

\subparagraph{Why the floor is essential, not parametric.}
A natural objection is that the freeze in the unfloored case
($\gamma_{\min} = 0$) might be a tuning issue resolved by a different
$(\delta_0, \rho)$ choice. Section~\ref{sec:setup} of
Chapter~\ref{chap:results} reports an empirical sweep across $25$
combinations of $\delta_0 \in \{0.01,\ldots,1.0\}\,f^{*}$ and
$\rho \in \{0.5,\ldots,0.99\}$ with $\gamma_{\min} = 0$ active and the
literal clip in place: only one configuration ($\delta_0 = 1.0\,f^{*}$,
$\rho = 0.99$) reaches a record gap below $10^{-3}$, and even there
the gap stalls two orders of magnitude above what the floored variant
achieves at the same parameter point. The floor is therefore not a
parameter fine-tune but a structural ingredient required for the
greedy-$\gamma$ deflection scheme to be operational on this problem
class --- a fact we attribute to ELM~LASSO's smooth-like behaviour
near the optimum, where consecutive subgradients align and the
parabola of eq.~\eqref{eq:gamma_parabola} pushes the unconstrained
minimiser below zero.

% ------------------------------------------------------------------
\subsection{Optimal deflection parameter}
\label{sec:dsm_optimal_gamma}

At each iteration on ELM LASSO, we choose $\gamma_i$ to minimize the deflected direction norm subject to a strictly positive lower bound $\gamma_{\min}$:
\begin{equation}\label{eq:gamma_opt}
\gamma_i \in \argmin_{\gamma \in [\gamma_{\min},1]} \norm{\gamma\, \mathbf{g}_i + (1-\gamma)\, \mathbf{d}_{i-1}}_2^2.
\end{equation}
Since the Polyak stepsize has denominator $\norm{\mathbf{d}_i}^2$, minimizing this maximizes step length and progress. The constraint $\gamma_i \ge \gamma_{\min}$ keeps the stepsize-restricted condition $\beta_i = \min(\beta,\gamma_i) \ge \gamma_{\min}$ operational (Section~\ref{sec:dsm_convergence_conditions}); we use $\gamma_{\min} = 0.05$ throughout and discuss the rationale in that section. The problem is quadratic in $\gamma$ with closed-form solution \cite{deflected-subgradient-method}:
\begin{equation}\label{eq:gamma_star}
\gamma^* = \frac{\norm{\mathbf{d}_{i-1}}_2^2 - \inner{\mathbf{g}_i}{\mathbf{d}_{i-1}}}{\norm{\mathbf{g}_i - \mathbf{d}_{i-1}}_2^2}, \quad \gamma_i = \min(1,\max(\gamma_{\min},\gamma^*)).
\end{equation}

\paragraph{Derivation of $\gamma^*$.}\mbox{} \\
We solve the unconstrained version of \eqref{eq:gamma_opt} over $\gamma \in \mathbb{R}$ and then project onto $[0,1]$.
Define $h(\gamma) := \norm{\gamma\,\mathbf{g}_i + (1-\gamma)\,\mathbf{d}_{i-1}}^2$.
Expanding by bilinearity of the inner product:
\begin{align}
h(\gamma)
&= \gamma^2\norm{\mathbf{g}_i}_2^2
  + 2\gamma(1-\gamma)\inner{\mathbf{g}_i}{\mathbf{d}_{i-1}}
  + (1-\gamma)^2\norm{\mathbf{d}_{i-1}}_2^2 \notag\\
&= \gamma^2\bigl(\norm{\mathbf{g}_i}_2^2
                - 2\inner{\mathbf{g}_i}{\mathbf{d}_{i-1}}
                + \norm{\mathbf{d}_{i-1}}_2^2\bigr)
  + 2\gamma\bigl(\inner{\mathbf{g}_i}{\mathbf{d}_{i-1}}
                - \norm{\mathbf{d}_{i-1}}_2^2\bigr)
  + \norm{\mathbf{d}_{i-1}}_2^2. \label{eq:gamma_expand}
\end{align}
Recognising the first factor as $\norm{\mathbf{g}_i - \mathbf{d}_{i-1}}^2$:
\begin{equation}
h(\gamma)
= \norm{\mathbf{g}_i - \mathbf{d}_{i-1}}_2^2\,\gamma^2
  + 2\bigl(\inner{\mathbf{g}_i}{\mathbf{d}_{i-1}} - \norm{\mathbf{d}_{i-1}}_2^2\bigr)\gamma
  + \norm{\mathbf{d}_{i-1}}_2^2.
\label{eq:gamma_parabola}
\end{equation}

\textbf{Case $\mathbf{g}_i = \mathbf{d}_{i-1}$:} the leading coefficient $\norm{\mathbf{g}_i - \mathbf{d}_{i-1}}_2^2 = 0$, so $h(\gamma) = \norm{\mathbf{d}_{i-1}}_2^2$ is constant on $[0,1]$; moreover, $\mathbf{d}_i = \gamma\mathbf{g}_i + (1-\gamma)\mathbf{d}_{i-1} = \mathbf{d}_{i-1}$ for all $\gamma$, so the deflected direction is identical regardless of the choice; we set $\gamma_i = 1$ by convention.

\textbf{Case $\mathbf{g}_i \neq \mathbf{d}_{i-1}$:} $h$ is a strictly convex parabola in $\gamma$ (positive leading coefficient). Its unique unconstrained minimizer satisfies $h'(\gamma^*) = 0$:
\[
h'(\gamma)
= 2\norm{\mathbf{g}_i - \mathbf{d}_{i-1}}_2^2\,\gamma
  + 2\bigl(\inner{\mathbf{g}_i}{\mathbf{d}_{i-1}} - \norm{\mathbf{d}_{i-1}}_2^2\bigr) = 0
\implies
\gamma^* = \frac{\norm{\mathbf{d}_{i-1}}_2^2 - \inner{\mathbf{g}_i}{\mathbf{d}_{i-1}}}
               {\norm{\mathbf{g}_i - \mathbf{d}_{i-1}}_2^2},
\]
which is exactly \eqref{eq:gamma_star}.
The constrained minimum over $[\gamma_{\min},1]$ is $\gamma_i = \min(1,\max(\gamma_{\min},\gamma^*))$: since the parabola opens upward, if $\gamma^* < \gamma_{\min}$ then $h$ is increasing on $[\gamma_{\min},1]$ and the minimum is attained at $\gamma = \gamma_{\min}$; if $\gamma^* > 1$ then $h$ is decreasing on the interval and the minimum is at $\gamma = 1$; if $\gamma^* \in [\gamma_{\min},1]$ the interior solution applies.

\paragraph{Interpretation.}\mbox{} \\
The denominator $\norm{\mathbf{g}_i - \mathbf{d}_{i-1}}_2^2$ measures the squared distance between the current subgradient and the previous direction: the larger it is, the more $\gamma^*$ shrinks and the more the deflected direction borrows from memory $\mathbf{d}_{i-1}$. The numerator $\norm{\mathbf{d}_{i-1}}_2^2 - \inner{\mathbf{g}_i}{\mathbf{d}_{i-1}}$ captures alignment: it is small (and $\gamma^*$ small) when $\mathbf{g}_i$ and $\mathbf{d}_{i-1}$ are well aligned, and large when they oppose each other. In the zig-zag scenario $\mathbf{g}_i \approx -\mathbf{d}_{i-1}$ with comparable magnitudes $\norm{\mathbf{g}_i}_2 \approx \norm{\mathbf{d}_{i-1}}_2$, the numerator is approximately $\norm{\mathbf{d}_{i-1}}_2^2 + \norm{\mathbf{d}_{i-1}}_2^2 = 2\norm{\mathbf{d}_{i-1}}_2^2$ while the denominator is approximately $\norm{2\mathbf{d}_{i-1}}_2^2 = 4\norm{\mathbf{d}_{i-1}}_2^2$, giving $\gamma^* \approx 1/2$: $\mathbf{d}_i$ splits evenly between the two conflicting directions, halving the oscillation.

% ------------------------------------------------------------------
\section{The target-level mechanism}
\label{sec:dsm_target_level}

The Polyak stepsize \cite{polyak-stepsize} requires knowledge of $f^*$. The target-level mechanism circumvents this by maintaining an adaptive estimate throughout the run.

The fundamental relationship for subgradient methods \cite{fundamental-relationships-sg} is:
\begin{equation}\label{eq:fundamental}
\norm{\mathbf{w}_{i+1}-\mathbf{w}^*}_2^2 \leq \norm{\mathbf{w}_i-\mathbf{w}^*}_2^2 - 2\alpha_i(f(\mathbf{w}_i)-f^*) + (\alpha_i)^2\norm{\mathbf{g}_i}_2^2.
\end{equation}
The ideal stepsize $\alpha_i^* = (f(\mathbf{w}_i)-f^*)/\norm{\mathbf{g}_i}_2^2$ \cite{polyak-stepsize} maximizes guaranteed reduction but requires $f^*$.

For ELM LASSO we replace $f^*$ with the target level $f_{\mathrm{ref}} - \delta$ \cite{target-level-mechanism}, where $f_{\mathrm{ref}}$ is the best function value found so far and $\delta > 0$ is the expected improvement; this is the same expression used in the Polyak-step numerator of Algorithm~\ref{alg:dsm}. When the target is accurate (i.e.\ $f_{\mathrm{ref}} - \delta \approx f^*$) the algorithm tracks the ideal Polyak step and achieves substantial progress; when the target is too optimistic (i.e.\ $f_{\mathrm{ref}} - \delta < f^*$), the numerator $f(\mathbf{w}_i) - (f_{\mathrm{ref}} - \delta) = (f(\mathbf{w}_i) - f^*) + (f^* - f_{\mathrm{ref}} + \delta)$ \emph{overestimates} $f(\mathbf{w}_i) - f^*$ and the resulting steps are larger than the ideal Polyak ones; the real failure mode is that the improvement condition $f(\mathbf{w}_{i+1}) \le f_{\mathrm{ref}} - \delta/2$ used by Algorithm~\ref{alg:dsm} can never be satisfied (since $f(\mathbf{w}_{i+1}) \ge f^* > f_{\mathrm{ref}} - \delta > f_{\mathrm{ref}} - \delta/2$), so no progress is registered and the iterates stagnate. We detect stalling by tracking the accumulated displacement $r = \sum \alpha_i \norm{\mathbf{d}_i}$: if $r > R$ (the patience threshold) without significant improvement, we contract $\delta \gets \delta \cdot \rho$ with $\rho < 1$ and retry. Repeated contractions drive $\delta \to 0$, so $f_{\mathrm{ref}} - \delta \to f_{\mathrm{ref}} \ge f^*$ and the target is restored above $f^*$.

% ------------------------------------------------------------------
\section{The complete algorithm}
\label{sec:dsm_algorithm}

Algorithm~\ref{alg:dsm} presents the full Deflected Subgradient Method with target level for ELM LASSO.

\begin{algorithm}[htbp]
\caption{Deflected Subgradient with Target Level (SGPTL) for ELM LASSO~\cite{deflected-subgradient-method}}
\label{alg:dsm}
\begin{algorithmic}[1]
\Require $f$ (ELM LASSO objective);\; $i_{\max}$;\; $\beta \in (0,1]$;\; $\delta_0 > 0$;\; $R > 0$;\; $\rho \in (0,1)$;\; $\gamma_{\min} \in (0,1]$
\State $\mathbf{w}_0 \gets (\mathbf{X}^T\mathbf{X})^{-1}\mathbf{X}^T\mathbf{y}$ \hfill\emph{(same OLS warm start as IRLS)}
\State $r \gets 0$;\; $\delta \gets \delta_0$;\; $f_{\mathrm{ref}} \gets \bar{f} \gets f(\mathbf{w}_0)$
\For{$i = 0, 1, \ldots, i_{\max}$}
\State Compute $\mathbf{g} \in \partial f(\mathbf{w}_i)$ \hfill\emph{(subgradient of LASSO)}
\If{$i = 0$} \hfill\emph{(no $\mathbf{d}_{-1}$ available)}
\State $\gamma_i \gets 1$;\; $\mathbf{d}_i \gets \mathbf{g}$
\Else
\State Compute $\gamma_i$ via \eqref{eq:gamma_star} \hfill\emph{(optimal deflection, clipped to $[\gamma_{\min},1]$)}
\State $\mathbf{d}_i \gets \gamma_i\, \mathbf{g} + (1-\gamma_i)\, \mathbf{d}_{i-1}$ \hfill\emph{(deflected direction)}
\EndIf
\State $\beta_i \gets \min(\beta,\gamma_i)$ \hfill\emph{(stepsize-restricted, $\beta=1$ default; see \S\ref{sec:dsm_convergence_conditions})}
\State $\alpha_i \gets \beta_i\,(f(\mathbf{w}_i) - (f_{\mathrm{ref}}-\delta))\,/\,\norm{\mathbf{d}_i}_2^2$ \hfill\emph{(Polyak step)}
\State $\mathbf{w}_{i+1} \gets \mathbf{w}_i - \alpha_i\, \mathbf{d}_i$ \hfill\emph{(update)}
\State $\bar{f} \gets \min(\bar{f},\; f(\mathbf{w}_{i+1}))$ \hfill\emph{(update record)}
\If{$f(\mathbf{w}_{i+1}) \leq f_{\mathrm{ref}} - \delta/2$} \hfill\emph{(sufficient descent)}
\State $f_{\mathrm{ref}} \gets \bar{f}$;\; $r \gets 0$
\ElsIf{$r > R$} \hfill\emph{(travel-distance patience exhausted)}
\State $\delta \gets \delta\,\rho$;\; $r \gets 0$
\Else
\State $r \gets r + \alpha_i\norm{\mathbf{d}_i}_2$
\EndIf
\EndFor
\State \Return $\mathbf{w}_{\mathrm{best}}$ (iterate achieving $\bar{f}$)
\end{algorithmic}
\end{algorithm}

\paragraph{Correspondence with the theory.}\mbox{} \\
Algorithm~\ref{alg:dsm} implements one-to-one the SGPTL rule of slide~14 of the course (Lessons~Optimization) and Lemma~3.8 of~\cite{dantonio2009deflected}. Lines~12--13 are the stepsize-restricted Polyak step ($\beta_i \le \gamma_i$, hypothesis of Theorem~\ref{thm:dsm_convergence}); lines~16--21 are the three-branch update of $(f_{\mathrm{ref}}, \delta, r)$ exactly as in Lemma~3.8. The single deviation from the slide's pseudocode is the clip on $\gamma_i$ at line~9: we project the greedy minimiser of~\eqref{eq:gamma_opt} onto $[\gamma_{\min}, 1]$ rather than $[0, 1]$, with $\gamma_{\min} = 0.05$. This is not a workaround but a structural choice that realises condition~(3.5) of~\cite{dantonio2009deflected} with $\beta^* = \gamma_{\min} > 0$: without the floor, the greedy choice~\eqref{eq:gamma_star} collapses to $0$ on near-aligned consecutive subgradients, $\beta_i = \min(\beta, \gamma_i) \to 0$ violates~(3.5), and the convergence hypothesis of Theorem~\ref{thm:dsm_convergence} fails. The floor is the minimal modification that keeps the theory applicable; we discuss the choice $\gamma_{\min} = 0.05$ in Section~\ref{sec:dsm_convergence_conditions}.

\paragraph{Initial-iteration convention.}\mbox{} \\
The first iteration $i=0$ is treated separately because~\eqref{eq:gamma_star} requires a previous direction $\mathbf{d}_{-1}$ that does not exist yet; setting $\gamma_0 = 1$ realises the convention $\mathbf{d}_{-1} = \mathbf{0}$ in eq.~\eqref{eq:deflected_dir}, so $\mathbf{d}_0 = \mathbf{g}_0$. Plugging $\mathbf{d}_{-1} = \mathbf{0}$ directly into~\eqref{eq:gamma_star} would yield $0/0$; the convention $\gamma_0 = 1$ is the same one used in the slide.

\paragraph{Numerical safeguards.}\mbox{} \\
Three pure floating-point safeguards complete the implementation; none changes the algorithmic behaviour, none reacts to algorithmic state (no $\delta$ contraction, no $\mathbf{d}$ reset), and none is in the pseudocode of Algorithm~\ref{alg:dsm}. \emph{(i)}~If $\norm{\mathbf{d}_i}^2 < 10^{-30}$ we exit the loop early: the deflected direction is numerically indistinguishable from zero, so $\alpha_i$ would divide by zero. This is a termination, not an alteration. \emph{(ii)}~If $f(\mathbf{w}_i) - (f_{\mathrm{ref}}-\delta) \le 0$, the Polyak numerator is non-positive and the step would point backwards. We set $\alpha_i = 0$ as prescribed by condition~(3.5) of~\cite{dantonio2009deflected}, and trigger the sufficient-descent update $f_{\mathrm{ref}} \gets \bar f$, $r \gets 0$ on $\mathbf{w}_i$ itself (since $f(\mathbf{w}_i) \le f_{\mathrm{ref}}-\delta \le f_{\mathrm{ref}}-\delta/2$): this is the theory-aligned response, not a $\delta$ contraction. \emph{(iii)}~If $\mathbf{w}_{i+1}$ contains a non-finite component (NaN/Inf), which can happen when $\norm{\mathbf{d}_i}^2$ is tiny but above the cutoff of~(i) and $\alpha_i$ overflows, we skip the iteration and proceed. None of $\delta$, $f_{\mathrm{ref}}$, $r$ is touched.

\paragraph{A note on what is \emph{not} in the algorithm.}\mbox{} \\
We deliberately do not include any non-theoretical safeguard such as an iteration-count fallback for $\delta$, an overshoot rescue snapping $\mathbf{w}$ to $\mathbf{w}_{\mathrm{best}}$ on large $f$-jumps, or extra contractions of $\delta$ outside line~19. Such mechanisms can rescue specific failure modes (the greedy collapse $\gamma_i \to \gamma_{\min}$ stalls the algorithm on ELM~LASSO well before the record converges to $f^*$, as we document in Chapter~\ref{chap:results}) but they have no support in the theory of~\cite{dantonio2009deflected} and would break the proof of Theorem~\ref{thm:dsm_convergence}; we prefer to expose the algorithm's empirical limitations honestly rather than mask them.

% ------------------------------------------------------------------
\subsection{Parameter calibration for ELM LASSO}
\label{sec:dsm_parameters}

SGPTL on ELM~LASSO has five parameters: the step modulation $\beta$, the deflection floor $\gamma_{\min}$, the initial target gap $\delta_0$, the contraction factor $\rho$, and the patience threshold $R$. We discuss them in order from the most theoretically constrained to the most empirically calibrated. The point worth stressing up front is that $f^*$ --- the optimum that the algorithm is trying to reach --- is \emph{not} an admissible input to any of these choices: tuning the optimiser against the quantity it is searching for is the optimisation analogue of training a model on its test set. All hyperparameters below are set from $f^*$-free quantities.

\paragraph{$\beta = 1$ (step modulation).}
\cite[(3.1), p.\,364]{dantonio2009deflected} admits any $\beta_k \in (0, 1]$ in the corrected Polyak stepsize. Slide~12 of the course (\cite{polyak-stepsize}) shows that, under the classical Polyak step $\alpha_k = \beta\,(f(\mathbf{w}_k) - f^*)/\norm{\mathbf{g}_k}^2$, the value $\beta = 1$ is exactly the minimiser of the one-step distance bound $\norm{\mathbf{w}_{k+1} - \mathbf{w}^*}^2$, i.e.\ the maximum-progress choice. The deflected variant in slide~15 inherits the same maximum-progress property, with the deflection coefficient $\gamma_i$ providing the dynamic damping; we therefore fix $\beta = 1$ and never sweep it. The effective per-iteration multiplier $\beta_i = \min(\beta, \gamma_i)$ in Algorithm~\ref{alg:dsm} is then driven entirely by $\gamma_i$ and by the floor $\gamma_{\min}$.

\paragraph{$\gamma_{\min} = 0.05$ (deflection floor).}
Condition~(3.5) of~\cite[p.\,365]{dantonio2009deflected} is a \emph{strictly necessary} hypothesis for the convergence of stepsize-restricted deflected methods: it requires $\beta_k \ge \beta^* > 0$ uniformly whenever $\lambda_k \ge 0$. With our $\beta_i = \min(1, \gamma_i)$, the uniform bound $\beta^*$ is realised exactly when $\gamma_i \ge \gamma_{\min} > 0$, in which case we may take $\beta^* = \gamma_{\min}$. The greedy choice~\eqref{eq:gamma_star} of $\gamma_i$ violates~(3.5) on ELM~LASSO (the closed-form minimiser collapses to $0$ on near-aligned consecutive subgradients, see~\S\ref{sec:dsm_convergence_conditions}), so the floor is the minimal modification that keeps the convergence theorem applicable. The value $\gamma_{\min} = 0.05$ is empirically calibrated: at this floor the clip activates only in degenerate steps (the unconstrained greedy $\gamma$ lies in $[\gamma_{\min}, 1]$ on a typical step), while the convergence rate $1/(\gamma_{\min}\sqrt k)$ of Theorem~\ref{thm:dsm_convergence} inflates by only $20\times$ relative to the unfloored Polyak rate. The sweep across $\gamma_{\min} \in \{0.05, 0.1, 0.2, 0.5\}$ in Chapter~\ref{chap:results} produces gaps consistently in $[10^{-7}, 10^{-5}]$, with $\gamma_{\min} = 0.05$ marginally best.

\paragraph{$\delta_0$ (initial target gap).}
The theory of \cite[Lemma~3.8]{dantonio2009deflected} requires only $\delta_0 \in (0, \infty)$; the convergence proof does not depend on the specific value, only on the contraction rule that drives $\delta_k \to 0$. The course slide~14 explicitly flags $\delta_0 > 0$ as ``\texttt{??}'' --- a free parameter. We therefore set $\delta_0 = c\,f(\mathbf{w}_0)$ where $\mathbf{w}_0$ is the OLS warm start (an $O(MH^2)$ pre-computation that does not see $f^*$) and $c = 0.1$. The scale $f(\mathbf{w}_0)$ is an upper bound on $f^*$ (since adding the regulariser strictly increases the value at a non-LASSO minimiser), and is the natural problem-dependent quantity for an order-of-magnitude estimate of the initial gap. The constant $c = 0.1$ is calibrated empirically in Section~\ref{sec:params}; the sweep across $c \in \{0.01, 0.05, 0.1, 0.5, 1\}$ produces gaps within a single order of magnitude, so the choice is not delicate. Section~\ref{sec:params} also compares this $f^*$-free choice against the diagnostic oracle $\delta_0 = 0.1\,f^*$ (reported for theoretical understanding only, never used in actual runs); the two are numerically indistinguishable on the tested instances.

\paragraph{$\rho = 0.7$ (threshold contraction).}
\cite[Lemma~3.8]{dantonio2009deflected} requires only $\mu \in (0, 1)$; the course slide~14 echoes the same open-interval constraint. No specific value is prescribed in any of the references we consulted (\cite{brannlund1993generalized,goffin1999convergence,bertsekas1999nonlinear}). Empirically we sweep $\rho \in \{0.3, 0.5, 0.7, 0.8, 0.9, 0.95\}$ in Section~\ref{sec:params} on the synthetic instance and on the two real-data ELM matrices; $\rho = 0.7$ is the value at most a factor-of-six away from the best on every tested instance, and is solidly inside the moderate range $[0.5, 0.8]$ where the target $f_{\mathrm{ref}} - \delta$ tracks $f^*$ at the right pace.

\paragraph{$R = 1$ (travel-distance patience).}
\cite[Lemma~3.8]{dantonio2009deflected} requires only $R > 0$; slide~14 flags it as ``\texttt{???}'', a free parameter not prescribed in any of \cite{brannlund1993generalized,goffin1999convergence}. The value $R$ controls the cadence of $\delta$-contractions: an $R$ much larger than the typical per-iteration step length $\alpha_i\norm{\mathbf{d}_i}$ leaves the $r_k > R$ branch inactive (the algorithm relies on the rare sufficient-descent firings to update $f^k_{\mathrm{ref}}$), an $R$ much smaller wastes the patience window. On ELM~LASSO with $\gamma_{\min} = 0.05$ and the OLS warm start, $\alpha_i\norm{\mathbf{d}_i}$ is of order $10^{-3}$, so we pick $R = 1$, which fires the contraction roughly every $10^3$ stalled iterations. An earlier draft of this report used $R = 10\sqrt{i_{\max}}\approx 894$ at $i_{\max}=8000$, three orders of magnitude above the natural scale, with the consequence that the $r > R$ branch never fired and the algorithm appeared to stall (see Section~\ref{sec:warm_cold} for the diagnosis). The choice $R = 1$ is theory-aligned --- only $R > 0$ is required --- and is the single calibration choice that the first submission got wrong.

% ------------------------------------------------------------------
\subsection{Starting point and patience-threshold calibration}
\label{sec:dsm_start}

Algorithm~\ref{alg:dsm} uses the same OLS warm start as IRLS, $\mathbf{w}_0 = (\mathbf{X}^\top\mathbf{X})^{-1}\mathbf{X}^\top\mathbf{y}$, so that both methods minimise the same problem from the same point and the empirical comparison is on equal footing. The travel-distance patience threshold is set to $R = 1$, a small absolute value chosen so the threshold is comparable to the per-iteration step length scale $\alpha_i\norm{\mathbf{d}_i} \sim 10^{-3}$ on ELM~LASSO with the OLS warm start and $\gamma_{\min} = 0.05$. The travel-distance branch then fires $\mathcal O(10)$--$\mathcal O(100)$ times per $8\,000$-iteration run, depending on $\rho$, contracting $\delta$ at a sustained pace as the theory expects (Section~\ref{sec:params} reports the rate). The choice $R = 1$ is within the theoretical framework of~\cite{dantonio2009deflected} (only $R > 0$ is required by Lemma~3.8) and is the only calibration choice that we make outside the slide's pseudocode; the rest of the parameters follow the slide directly.

\paragraph{Implementation history and rationale.}
The first submission of this project used additional mechanisms on top of Algorithm~\ref{alg:dsm}: (a)~an iteration-count $\delta$-contraction trigger (auxiliary counter $c$ tracking iterations without improvement of $\bar f$, with the rule ``if $c \ge R_{\mathrm{iter}}$ then $\delta \gets \rho\delta$ and $\mathbf{d}_{i-1} \gets \mathbf{0}$''); (b)~an overshoot rescue that snapped $\mathbf{w}$ back to $\mathbf{w}_{\mathrm{best}}$ on large Polyak-step blow-ups and contracted $\delta$; and (c)~the default $R = 10\sqrt{i_{\max}} \approx 894$ for the travel-distance threshold, which on ELM~LASSO was orders of magnitude larger than any realistic accumulated travel between resets, so the $r > R$ branch effectively never fired and the iteration-count fallback was the sole driver of $\delta$-contractions. The instructor's feedback flagged this construction as inconsistent with the theory of~\cite{dantonio2009deflected}: the proof of $\sum_k s_k = \infty$ in Lemma~3.8 requires every $\delta$-contraction to be preceded by at least $R$ of accumulated travel via the $r_k > R$ trigger, which the iteration-count fallback breaks. Theorem~\ref{thm:dsm_convergence}, which cites~\cite[Theorem~3.7]{dantonio2009deflected} applied to Lemma~3.8 verbatim, was therefore not fully applicable to the algorithm we were actually running.

The diagnosis after the feedback was twofold. First, mechanisms (a) and (b) are workarounds that hide bugs rather than fix them, and the proof technique of~\cite[Lemma~3.8]{dantonio2009deflected} does not extend to them; they should not be in the algorithm. Second, mechanism (c) --- a single integer in our defaults --- was the actual reason the theory-pure algorithm appeared to ``not work'' in the first submission: with $R$ chosen at the right scale, the $r > R$ branch of Lemma~3.8 fires the way the theory needs it to, $\delta$ contracts at a sustained pace, and the algorithm converges within the theoretical envelope of Theorem~\ref{thm:dsm_convergence}. The starting point, on the other hand, turns out to be approximately neutral: with $R = 1$ both the OLS warm start and a cold start $\mathbf{w}_0 = \mathbf{0}$ produce sublinear descent, with the warm start slightly behind on some instances and slightly ahead on others (Section~\ref{sec:warm_cold} documents the side-by-side comparison). The second submission therefore: removes (a) and (b); sets the default to $R = 1$; keeps the OLS warm start for both algorithms to retain the apples-to-apples comparison. The empirical cost of switching to the theory-pure rule is moderate (gap $\sim 10^{-5}$ at $8\,000$ iterations on the convergence instance, instead of the $\sim 10^{-8}$ that the workarounds produced); we view this as the honest baseline of the algorithm and discuss the gap to IRLS in the results and conclusions chapters.


% ------------------------------------------------------------------
\section{Application to the ELM LASSO problem}
\label{sec:dsm_application}

\subsection{Subgradient computation for ELM LASSO}
\label{sec:dsm_subgrad}

For the ELM LASSO objective, decompose $f = g + h$ where $g(\mathbf{w}) = \frac{1}{2}\norm{\mathbf{X}\mathbf{w}-\mathbf{y}}_2^2$ (smooth) and $h(\mathbf{w}) = \lambda_{\text{LASSO}}\norm{\mathbf{w}}_1$ (nonsmooth).

By the sum rule for subdifferentials \cite{subdifferential-sum-rule}:
\[ \partial f(\mathbf{w}) = \nabla g(\mathbf{w}) + \lambda_{\text{LASSO}}\,\partial\norm{\mathbf{w}}_1, \]
where $\nabla g(\mathbf{w}) = \mathbf{X}^T(\mathbf{X}\mathbf{w} - \mathbf{y})$ and the subdifferential of the $\ell_1$ norm is component-wise (Section~\ref{sec:intro_optimality}). A subgradient of $f$ at $\mathbf{w}$ is:
\begin{equation}\label{eq:dsm_subgrad}
  \mathbf{g} = \mathbf{X}^T(\mathbf{X}\mathbf{w} - \mathbf{y}) + \lambda_{\text{LASSO}}\, \mathbf{s}, \quad s_i \in \partial|w_i|.
\end{equation}
At a component where $w_i \ne 0$ the subdifferential is a singleton and $s_i = \operatorname{sign}(w_i)$ is forced; at $w_i = 0$ the subdifferential is the whole interval $[-1, +1]$ and we must pick a value. Writing $\mathbf{q} = \mathbf{X}^\top(\mathbf{X}\mathbf{w}-\mathbf{y})$ for the smooth gradient, the minimum-norm element of $\partial f$ at an index with $w_i = 0$ is obtained by minimising $(q_i + \lambda_{\text{LASSO}}\,s_i)^2$ over $s_i \in [-1, +1]$, which gives $s_i^{\star} = -q_i/\lambda_{\text{LASSO}}$ when $|q_i| \le \lambda_{\text{LASSO}}$ (so that $g_i = 0$) and $s_i^{\star} = -\operatorname{sign}(q_i)$ otherwise.

Our implementation takes $s_i = 0$ regardless, via \texttt{numpy.sign(0)\,=\,0}. This value is always admissible --- it lies in $[-1, +1]$ --- so $\mathbf{g}$ is a valid subgradient and the convergence guarantee of Theorem~\ref{thm:dsm_convergence} applies as written. It is \emph{not} the minimum-norm element of $\partial f$ at $\mathbf{w}$: the unconstrained minimum of $(q_i + \lambda_{\text{LASSO}} s_i)^2$ over $s_i$ is $s_i^{\star} = -q_i/\lambda_{\text{LASSO}}$, which yields $g_i^{\star} = 0$ when $|q_i| \le \lambda_{\text{LASSO}}$ and $s_i^{\star} = -\operatorname{sign}(q_i)$ with $|g_i^{\star}| = |q_i| - \lambda_{\text{LASSO}}$ otherwise; our choice $s_i = 0$ matches neither (except trivially when $q_i = 0$). The discrepancy in $|g_i|$ is at most $\lambda_{\text{LASSO}}$ per affected component, so it inflates the constant $L$ in the convergence rate but not the $O(L^2/\varepsilon^2)$ order. In practice the issue is essentially theoretical for SGPTL on the OLS warm start: the iterate $w_i$ equals zero exactly only at indices that were zero in $\mathbf{w}_0$, an event of zero Lebesgue measure under the Polyak-step update once the algorithm leaves the initial point.

% ------------------------------------------------------------------
\subsection{Convergence analysis and hypothesis verification for ELM LASSO}
\label{sec:dsm_convergence}

\begin{theorem}[Convergence of SGPTL, adapted from \cite{target-level-mechanism,deflected-subgradient-method,dantonio2009deflected}]
\label{thm:dsm_convergence}
Let $f : \R^H \to \R$ be convex with $f^* = \min_{\mathbf{w}} f(\mathbf{w}) > -\infty$ attained at some $\mathbf{w}^*$, and let $\{\mathbf{w}_i\}$ be the sequence generated by Algorithm~\ref{alg:dsm} (the target-level deflected subgradient method we actually run). Assume the subgradients encountered along the run are uniformly bounded, $\norm{\mathbf{g}_i}_2 \leq L$, and the stepsize-restricted rule $\beta_i \le \gamma_i$ holds with $\gamma_i \ge \gamma_{\min} > 0$. Then the record values $\bar{f}_i = \min_{h \le i} f(\mathbf{w}_h)$ converge to $f^*$, and once the target-level estimate has sharpened to $f^*$ (i.e.\ $\delta_k \to 0$ and $f^k_{\mathrm{ref}} \to f^*$, which Lemma~3.8 of \cite{dantonio2009deflected} guarantees) the asymptotic rate matches the classical Polyak bound \cite{polyak-rate-frangioni,bertsekas1999nonlinear} inflated by the deflection floor:
\begin{equation}
\label{eq:def-rate-deflected}
\bar{f}_k - f^* \;\leq\; \frac{L\,\norm{\mathbf{w}_0 - \mathbf{w}^*}_2}{\gamma_{\min}\sqrt{k}}\,,
\end{equation}
so the worst-case complexity to attain $\bar f_k - f^* \le \varepsilon$ is $O(L^2/(\gamma_{\min}^2\varepsilon^2))$.
\end{theorem}

\begin{proof}
We verify the hypotheses of \cite[Theorem~3.7]{dantonio2009deflected} on Algorithm~\ref{alg:dsm} and read off the conclusions. The notation in \cite{dantonio2009deflected} maps to ours via $\sigma_k \leftrightarrow 0$ (exact oracle), $\alpha_k \leftrightarrow \gamma_i$, $\beta_k \leftrightarrow \beta_i$, $\nu_k \leftrightarrow \alpha_i$, $\mu \leftrightarrow \rho$, $X = \R^H$ (no constraint), and $\lambda_k \leftrightarrow f(\mathbf{w}_i) - (f^k_{\mathrm{ref}} - \delta_k)$.

\emph{(a)~Deflection-consistency condition~(2.13).} In \cite{dantonio2009deflected} condition~(2.13) of Lemma~2.4 reads $d_k = \tilde d_k \Rightarrow v_{k+1} = \tilde d_k$, where $\tilde d_k$ is the unprojected deflected direction and $\hat d_k = -P_{T_k}(-\tilde d_k)$ the projected one. This is the hypothesis under which the ``crucial'' inequality (2.12) holds, which the stepsize-restricted analysis relies on. With $X = \R^H$ the tangent cone is $T_k = \R^H$ everywhere, so $\hat d_k = \tilde d_k$ and necessarily $d_k = \tilde d_k$ at every iteration; the algorithm sets $v_{k+1}$ to $d_k = \tilde d_k$ (Algorithm~\ref{alg:dsm} carries no projection step), so~(2.13) is satisfied trivially.

\emph{(b)~Stepsize condition (3.5)}: $\lambda_k \ge 0 \Rightarrow \alpha_k \ge \beta_k \ge \beta^* > 0$, $\lambda_k < 0 \Rightarrow \alpha_k = 0$. When $\lambda_k > 0$, line~12 of Algorithm~\ref{alg:dsm} sets $\beta_i = \min(1,\gamma_i) \ge \gamma_{\min}$ via the floor on $\gamma_i$, so we may take $\beta^* = \gamma_{\min}$; the inequality $\alpha_k \ge \beta_k$ in (3.5) is $\gamma_i \ge \beta_i$, which is exactly the stepsize-restricted rule. When $\lambda_k \le 0$ the numerical safeguard of Section~\ref{sec:dsm_algorithm} skips the step --- equivalent to $\alpha_i = 0$ --- and triggers the sufficient-descent update $f^k_{\mathrm{ref}} \gets \bar f$, $r \gets 0$ on $\mathbf{w}_i$ itself. The skip realises (3.5) verbatim; the additional $f^k_{\mathrm{ref}}$ update is theory-aligned because $\lambda_k \le 0$ means $f(\mathbf{w}_i) \le f^k_{\mathrm{ref}} - \delta_k \le f^k_{\mathrm{ref}} - \delta_k/2$, which is the sufficient-descent condition of \cite[Lemma~3.8]{dantonio2009deflected} evaluated at $\mathbf{w}_i$ rather than at $\mathbf{w}_{i+1}$.

\emph{(c)~Target-level threshold condition~(3.26)}: requires \emph{either} $f^\infty_{\mathrm{ref}} = f^* = -\infty$ \emph{or} $\liminf_k \delta_k = 0$ together with the series condition~(3.25), $\sum_k \lambda_k/\norm{\mathbf{d}_k}^2 = \infty$ (equivalent to $\sum_k s_k = \infty$ via~(3.28) of \cite{dantonio2009deflected}, where $s_k = \norm{\mathbf{w}_{k+1}-\mathbf{w}_k} = \beta_k \lambda_k / \norm{\mathbf{d}_k}$). Since on ELM~LASSO $f^* > -\infty$, the first branch is excluded and we need the second. \cite[Lemma~3.8]{dantonio2009deflected} provides a concrete update strategy for $(f^k_{\mathrm{ref}}, \delta_k)$ --- the sufficient-descent test on $f_{k+1} \le f^k_{\mathrm{ref}} - \delta_k/2$, the target-infeasibility test on $r_k > R$, and the accumulate-$r$ counter --- which attains both $\delta_k \to 0$ and $\sum_k s_k = \infty$. Lines~16--21 of Algorithm~\ref{alg:dsm} implement that strategy verbatim under the identification $\mu \leftrightarrow \rho$, so condition~(3.26) holds.

\emph{Subgradients uniformly bounded.} The convex objective $f$ of ELM~LASSO is coercive and continuous; the iterates $\{\mathbf{w}_i\}$ remain in a compact sublevel set $S_{c_0}$ for some $c_0 \ge f(\mathbf{w}_0)$ (verification below), and $\partial f$ is uniformly bounded on $S_{c_0}$ by Lipschitz continuity of convex functions on compact sets \cite{lipschitz-continuity-frangioni}.

Conditions (2.13), (3.5), (3.26) being verified with $\sigma^* = 0$, \cite[Theorem~3.7]{dantonio2009deflected} gives $f^\infty_{\mathrm{ref}} \le f^* + \sigma^* = f^*$, so $\bar f_i \to f^*$. The asymptotic rate~\eqref{eq:def-rate-deflected} is then the classical Polyak rate \cite[Theorem~7.17]{bertsekas1999nonlinear}, \cite{polyak-rate-frangioni} applied to the regime where the target-level estimate has sharpened to $f^*$ (Lemma~3.8 of \cite{dantonio2009deflected} guarantees $\delta_k \to 0$, so eventually the step reduces to the standard Polyak step with $f^*$); the factor $1/\gamma_{\min}$ is the explicit price of the deflection floor, which contracts the effective stepsize by $\beta_i = \min(1, \gamma_i) \ge \gamma_{\min}$ at every iteration. On ELM~LASSO with $\gamma_{\min} = 0.05$ this inflates the iteration count to reach a given $\varepsilon$ by $400\times$ relative to the unfloored Polyak rate, but the $O(L^2/\varepsilon^2)$ order is preserved.
\end{proof}

\paragraph{Verification of theorem assumptions for ELM LASSO.}

\begin{enumerate}
\item \textbf{Convexity of } $f(\mathbf{w}) = \frac{1}{2}\norm{\mathbf{X}\mathbf{w}-\mathbf{y}}_2^2 + \lambda_{\text{LASSO}}\norm{\mathbf{w}}_1$.

The smooth part $g$ has Hessian $\nabla^2 g = \mathbf{X}^T\mathbf{X} \succeq 0$ (positive semidefinite), so $g$ is convex. The $\ell_1$ norm is convex. The sum of convex functions is convex \cite{convex-functions-frangioni}, so $f$ is convex.

We assume $\mathbf{X}$ has full column rank, as commonly observed in ELM practice when $M \ge H$ \cite{HUANG2006489}; under this assumption $\nabla^2 g \succ 0$, so $g$ is $\mu$-\emph{strongly convex} with modulus $\mu = \lambda_{\min}(\mathbf{X}^T\mathbf{X}) > 0$, and adding the convex term $h$ preserves the property \cite{strongly-convex-frangioni,convex-functions-frangioni}, so $f = g + h$ is $\mu$-strongly convex as well. In principle this structure could be exploited by methods designed for strongly-convex nonsmooth problems, which improve the worst-case complexity to $O(L^2/(\mu\varepsilon))$ \cite{strongly-convex-frangioni}. However, Algorithm~\ref{alg:dsm} does not incorporate any such mechanism: deflection only dampens oscillations, it does not introduce the kind of acceleration or proximal step needed to exploit $\mu$, so the worst-case guarantee for our SGPTL on ELM LASSO remains the standard $O(L^2/\varepsilon^2)$.

\item \textbf{Bounded subgradients for ELM LASSO.}

The $\ell_1$ component has $\norm{\lambda_{\text{LASSO}} \mathbf{s}}_2 \leq \lambda_{\text{LASSO}}\sqrt{H}$ (since each $|s_i| \leq 1$). The smooth component $\mathbf{X}^T(\mathbf{X}\mathbf{w}-\mathbf{y})$ grows unboundedly as $\norm{\mathbf{w}} \to \infty$, so $f$ is not globally Lipschitz on $\mathbb{R}^H$.

However, $f$ is coercive ($f(\mathbf{w}) \to \infty$ as $\|\mathbf{w}\| \to \infty$), so every sublevel set $S_c = \{\mathbf{w} : f(\mathbf{w}) \le c\}$ is compact. Subgradient methods do not guarantee monotone decrease of $f(\mathbf{w}_i)$, so individual iterates may temporarily leave $S_{f(\mathbf{w}_0)}$; what stays inside is the best iterate, since $\bar{f}_i = \min_{h \le i} f(\mathbf{w}_h) \le f(\mathbf{w}_0)$. In practice the target-level mechanism keeps the steps bounded: the numerator $f(\mathbf{w}_i) - (f_{\mathrm{ref}} - \delta)$ stays of the same order of magnitude as $f(\mathbf{w}_i) - f^*$, so the iterates do not drift far from $\mathbf{w}_{\mathrm{best}}$ and remain inside an enlarged compact set $S_{c_0}$ for some $c_0 \ge f(\mathbf{w}_0)$ depending on the problem data. By Lipschitz continuity of $f$ on this compact set \cite{lipschitz-continuity-frangioni}, all subgradients encountered by the algorithm are uniformly bounded by some $L$ depending on the problem data and the initial point.

\end{enumerate}

\paragraph{Complexity context: why SGPTL for ELM LASSO.}\mbox{} \\
ELM LASSO is nonsmooth and convex. The intrinsic lower bound is $\Omega(L^2/\varepsilon^2)$ for any first-order method on nonsmooth convex problems \cite{nonsmooth-lower-bounds-frangioni,bubeck2015convex}. Algorithm~\ref{alg:dsm} achieves this bound, making it theoretically optimal. Deflection does not improve worst-case complexity; its benefit is reducing oscillations and improving practical constants in the $O$ notation.

% ------------------------------------------------------------------
\section{Expected practical behavior on the ELM LASSO problem}
\label{sec:dsm_properties}

SGPTL has three features that the experimental chapter will need to take into account. Convergence is \emph{non-monotone}: $f(\mathbf{w}_{i+1}) > f(\mathbf{w}_i)$ occurs frequently and only the record $\bar{f}_i = \min_{h \le i} f(\mathbf{w}_h)$ decreases monotonically to $f^*$, so semilog plots should display $\log(\bar{f}_i - f^*)$ and the returned iterate is $\mathbf{w}_{\mathrm{best}}$, not $\mathbf{w}_{i_{\max}}$. The rate is sublinear, $O(L^2/\varepsilon^2)$ from Theorem~\ref{thm:dsm_convergence}, so doubling the target precision roughly quadruples the iteration count and accuracies below $10^{-6}$ are not a realistic goal for this method. The per-iteration cost is $O(MH)$ (two matrix-vector products with $\mathbf{X}$), against the $O(H^3)$ Cholesky solve of IRLS, so DSM is competitive on CPU time only when $H$ is large and the accuracy target is moderate.

Two further points will matter at tuning time. There is no reliable subgradient-based stopping test: unlike the smooth case, where $\norm{\nabla f(\mathbf{w})} = 0$ signals optimality, here $\mathbf{w}^*$ satisfies only $0 \in \partial f(\mathbf{w}^*)$ and other elements of the subdifferential at $\mathbf{w}^*$ can have non-zero norm, so $\norm{\mathbf{g}_i}$ remaining large is uninformative; we therefore stop on the iteration budget $i_{\max}$ or when $\bar{f}_i$ fails to improve appreciably over a fixed window. Sparsity is also different from IRLS, although the difference is more nuanced than it may look at first sight: neither method produces exact zeros without a post-termination threshold. IRLS components corresponding to true zeros are pinned around the safety floor $\varepsilon_{\mathrm{thr}}$ in eq.~(\ref{eq:irls_weights}) (which we set to $10^{-8}$), while SGPTL leaves them at the residual-rate level $\sim L/\sqrt{i}$, typically $10^{-3}$--$10^{-4}$ at the budgets we run. Both methods therefore require thresholding to report sparsity; using the same threshold on both sides is what makes the comparison fair, and is the convention we adopt in Section~\ref{sec:quality}. Finally, the regularization strength $\lambda_{\text{LASSO}}$ enters the subgradient directly through $\mathbf{s}$, so larger values amplify oscillations and slow convergence; the target-level mechanism adapts automatically, but the iteration count is expected to grow with $\lambda_{\text{LASSO}}$.
